\section{Experiments}
\label{sec:experiments}
\newcommand{\emptycell}{\textemdash}

We evaluate our method on Stable Diffusion 3.5 Medium using four experiment groups: I2P~\cite{Schramowski_2023_CVPR} with MS-COCO~\cite{lin2015microsoftcococommonobjects} for NSFW suppression and general generation preservation, IP character erasure, artistic style erasure, and object erasure.
The IP and object experiments include both single-concept and multiple-concept erasure settings.
We compare with UCE~\cite{Gandikota_2024_WACV}, ESD~\cite{Gandikota_2023_ICCV}, EraseAnything~\cite{pmlr-v267-gao25j}, RAZOR~\cite{Ranjan_2026_CVPR}, and CASteer~\cite{gaintseva2026casteercrossattentionsteeringcontrollable}, using their SD3.5 adaptations in our benchmark.

\noindent\textbf{Implementation details.}
We optimize both prompt types with AdamW at a learning rate of $10^{-3}$ for up to five epochs.
The SupCon weight and temperature are set to $\lambda=1$ and $T=0.1$, respectively, and the detection threshold is $\tau_c=0.2$.
Images are generated at $512\times512$ resolution using the model's native flow matching Euler sampler with 40 steps and guidance scale $s=4.5$.
All methods share the same evaluation prompts, seeds, and base sampling settings.

\noindent\textbf{Evaluation metrics.}
For NSFW suppression, we report the fraction of images flagged by NudeNet.
On MS-COCO, we compute CLIP score (CS) using CLIP ViT-B/16~\cite{pmlr-v139-radford21a} as the mean cosine similarity between image and original prompt embeddings.
FID~\cite{NIPS2017_8a1d6947} measures distribution change relative to images generated by the original model under the same sampling protocol.
For IP and style, we report FID separately for each target and non-target concept.
Higher target FID indicates greater distribution change, while lower non-target FID indicates stronger preservation.
Qwen3-VL-2B-Instruct~\cite{bai2025qwen3vltechnicalreport} measures character and style recognition by assigning each image to one of three candidate concepts or \texttt{NONE}, using only the image.
Object recognition is evaluated with ResNet-50~\cite{He_2016_CVPR} using ImageNet-1K V2 weights and Top-1/Top-3 predictions over all 1,000 classes.
Lower target recognition and higher non-target recognition indicate stronger erasure and preservation, respectively.

\subsection{NSFW erasure and general generation preservation}
\label{sec:exp_nsfw}

We evaluate NSFW suppression on the 931 prompts in the I2P sexual subset, generating one image per prompt with its provided seed.
NudeNet flags an image when at least one of eight exposure categories is detected with confidence at least $0.6$.
For general generation preservation, we use the first 1,000 captions in the fixed MS-COCO caption list, generating one image per caption.
Both evaluations use the same nudity suppression prompt and detection settings.
\Cref{tab:nsfw_preservation} reports the I2P detection rate alongside MS-COCO CS and FID.

\begin{table}[t]
    \centering
    \small
    \setlength{\tabcolsep}{4pt}
    \caption{NSFW suppression on I2P and general generation preservation on MS-COCO. NudeNet detection rates are reported as percentages.}
    \label{tab:nsfw_preservation}
    \begin{tabular}{@{}lccc@{}}
        \toprule
        Method & I2P & \multicolumn{2}{c}{MS-COCO} \\
        \cmidrule(lr){2-2}\cmidrule(lr){3-4}
        & NudeNet $\downarrow$ & CS $\uparrow$ & FID $\downarrow$ \\
        \midrule
        SD3.5 Medium & \emptycell & \emptycell & \emptycell \\
        UCE & \emptycell & \emptycell & \emptycell \\
        ESD & \emptycell & \emptycell & \emptycell \\
        EraseAnything & \emptycell & \emptycell & \emptycell \\
        RAZOR & \emptycell & \emptycell & \emptycell \\
        CASteer & \emptycell & \emptycell & \emptycell \\
        Ours & \emptycell & \emptycell & \emptycell \\
        \bottomrule
    \end{tabular}
\end{table}

\subsection{IP character erasure}
\label{sec:exp_ip}

We evaluate Pikachu, Snoopy, and Mickey using 80 fixed prompt templates per character and seeds 0--3, yielding 960 images per setting.
The first setting erases Pikachu, while the second enables erasure of both Pikachu and Snoopy in the same model.
All three characters are requested in both settings to assess target removal and preservation of the remaining characters.
\Cref{tab:ip_erasure} reports FID and VLM recognition separately for each character.

\begin{table*}[t]
    \centering
    \small
    \setlength{\tabcolsep}{2pt}
    \caption{IP character erasure under single-target and two-target settings. Each character has FID and VLM recognition (\%) columns. Target concepts are italicized.}
    \label{tab:ip_erasure}
    \begin{subtable}[t]{0.49\textwidth}
        \centering
        \caption{Erasing Pikachu}
        \label{tab:ip_single}
        \begin{tabular}{@{}lcccccc@{}}
            \toprule
            Method & \multicolumn{2}{c}{\textit{Pikachu}} & \multicolumn{2}{c}{Snoopy} & \multicolumn{2}{c}{Mickey} \\
            \cmidrule(lr){2-3}\cmidrule(lr){4-5}\cmidrule(lr){6-7}
            & FID $\uparrow$ & VLM $\downarrow$ & FID $\downarrow$ & VLM $\uparrow$ & FID $\downarrow$ & VLM $\uparrow$ \\
            \midrule
            SD3.5 Medium & \emptycell & \emptycell & \emptycell & \emptycell & \emptycell & \emptycell \\
            UCE & \emptycell & \emptycell & \emptycell & \emptycell & \emptycell & \emptycell \\
            ESD & \emptycell & \emptycell & \emptycell & \emptycell & \emptycell & \emptycell \\
            EraseAnything & \emptycell & \emptycell & \emptycell & \emptycell & \emptycell & \emptycell \\
            RAZOR & \emptycell & \emptycell & \emptycell & \emptycell & \emptycell & \emptycell \\
            CASteer & \emptycell & \emptycell & \emptycell & \emptycell & \emptycell & \emptycell \\
            Ours & \emptycell & \emptycell & \emptycell & \emptycell & \emptycell & \emptycell \\
            \bottomrule
        \end{tabular}
    \end{subtable}
    \hfill
    \begin{subtable}[t]{0.49\textwidth}
        \centering
        \caption{Erasing Pikachu and Snoopy}
        \label{tab:ip_multi}
        \begin{tabular}{@{}lcccccc@{}}
            \toprule
            Method & \multicolumn{2}{c}{\textit{Pikachu}} & \multicolumn{2}{c}{\textit{Snoopy}} & \multicolumn{2}{c}{Mickey} \\
            \cmidrule(lr){2-3}\cmidrule(lr){4-5}\cmidrule(lr){6-7}
            & FID $\uparrow$ & VLM $\downarrow$ & FID $\uparrow$ & VLM $\downarrow$ & FID $\downarrow$ & VLM $\uparrow$ \\
            \midrule
            SD3.5 Medium & \emptycell & \emptycell & \emptycell & \emptycell & \emptycell & \emptycell \\
            UCE & \emptycell & \emptycell & \emptycell & \emptycell & \emptycell & \emptycell \\
            ESD & \emptycell & \emptycell & \emptycell & \emptycell & \emptycell & \emptycell \\
            EraseAnything & \emptycell & \emptycell & \emptycell & \emptycell & \emptycell & \emptycell \\
            RAZOR & \emptycell & \emptycell & \emptycell & \emptycell & \emptycell & \emptycell \\
            CASteer & \emptycell & \emptycell & \emptycell & \emptycell & \emptycell & \emptycell \\
            Ours & \emptycell & \emptycell & \emptycell & \emptycell & \emptycell & \emptycell \\
            \bottomrule
        \end{tabular}
    \end{subtable}
\end{table*}

\subsection{Artistic style erasure}
\label{sec:exp_style}

We conduct three independent experiments, erasing Van Gogh, Monet, and Picasso, respectively.
Each setting is evaluated on all three styles using 30 painting templates per style and seeds 0--2, yielding 270 images.
The target style measures erasure, while the remaining two styles measure preservation.
\Cref{tab:style_erasure} reports FID and VLM style recognition for each style in each setting.

\begin{table}[t]
    \centering
    \small
    \setlength{\tabcolsep}{2pt}
    \caption{Artistic style erasure, with a separate setting for each target artist. Each style has FID and VLM style recognition (\%) columns. Target styles are italicized.}
    \label{tab:style_erasure}
    \begin{subtable}{\columnwidth}
        \centering
        \caption{Erasing Van Gogh}
        \label{tab:style_van_gogh}
        \begin{tabular}{@{}lcccccc@{}}
            \toprule
            Method & \multicolumn{2}{c}{\textit{Van Gogh}} & \multicolumn{2}{c}{Monet} & \multicolumn{2}{c}{Picasso} \\
            \cmidrule(lr){2-3}\cmidrule(lr){4-5}\cmidrule(lr){6-7}
            & FID $\uparrow$ & VLM $\downarrow$ & FID $\downarrow$ & VLM $\uparrow$ & FID $\downarrow$ & VLM $\uparrow$ \\
            \midrule
            SD3.5 Medium & \emptycell & \emptycell & \emptycell & \emptycell & \emptycell & \emptycell \\
            UCE & \emptycell & \emptycell & \emptycell & \emptycell & \emptycell & \emptycell \\
            ESD & \emptycell & \emptycell & \emptycell & \emptycell & \emptycell & \emptycell \\
            EraseAnything & \emptycell & \emptycell & \emptycell & \emptycell & \emptycell & \emptycell \\
            RAZOR & \emptycell & \emptycell & \emptycell & \emptycell & \emptycell & \emptycell \\
            CASteer & \emptycell & \emptycell & \emptycell & \emptycell & \emptycell & \emptycell \\
            Ours & \emptycell & \emptycell & \emptycell & \emptycell & \emptycell & \emptycell \\
            \bottomrule
        \end{tabular}
    \end{subtable}
    \par\smallskip
    \begin{subtable}{\columnwidth}
        \centering
        \caption{Erasing Monet}
        \label{tab:style_monet}
        \begin{tabular}{@{}lcccccc@{}}
            \toprule
            Method & \multicolumn{2}{c}{Van Gogh} & \multicolumn{2}{c}{\textit{Monet}} & \multicolumn{2}{c}{Picasso} \\
            \cmidrule(lr){2-3}\cmidrule(lr){4-5}\cmidrule(lr){6-7}
            & FID $\downarrow$ & VLM $\uparrow$ & FID $\uparrow$ & VLM $\downarrow$ & FID $\downarrow$ & VLM $\uparrow$ \\
            \midrule
            SD3.5 Medium & \emptycell & \emptycell & \emptycell & \emptycell & \emptycell & \emptycell \\
            UCE & \emptycell & \emptycell & \emptycell & \emptycell & \emptycell & \emptycell \\
            ESD & \emptycell & \emptycell & \emptycell & \emptycell & \emptycell & \emptycell \\
            EraseAnything & \emptycell & \emptycell & \emptycell & \emptycell & \emptycell & \emptycell \\
            RAZOR & \emptycell & \emptycell & \emptycell & \emptycell & \emptycell & \emptycell \\
            CASteer & \emptycell & \emptycell & \emptycell & \emptycell & \emptycell & \emptycell \\
            Ours & \emptycell & \emptycell & \emptycell & \emptycell & \emptycell & \emptycell \\
            \bottomrule
        \end{tabular}
    \end{subtable}
    \par\smallskip
    \begin{subtable}{\columnwidth}
        \centering
        \caption{Erasing Picasso}
        \label{tab:style_picasso}
        \begin{tabular}{@{}lcccccc@{}}
            \toprule
            Method & \multicolumn{2}{c}{Van Gogh} & \multicolumn{2}{c}{Monet} & \multicolumn{2}{c}{\textit{Picasso}} \\
            \cmidrule(lr){2-3}\cmidrule(lr){4-5}\cmidrule(lr){6-7}
            & FID $\downarrow$ & VLM $\uparrow$ & FID $\downarrow$ & VLM $\uparrow$ & FID $\uparrow$ & VLM $\downarrow$ \\
            \midrule
            SD3.5 Medium & \emptycell & \emptycell & \emptycell & \emptycell & \emptycell & \emptycell \\
            UCE & \emptycell & \emptycell & \emptycell & \emptycell & \emptycell & \emptycell \\
            ESD & \emptycell & \emptycell & \emptycell & \emptycell & \emptycell & \emptycell \\
            EraseAnything & \emptycell & \emptycell & \emptycell & \emptycell & \emptycell & \emptycell \\
            RAZOR & \emptycell & \emptycell & \emptycell & \emptycell & \emptycell & \emptycell \\
            CASteer & \emptycell & \emptycell & \emptycell & \emptycell & \emptycell & \emptycell \\
            Ours & \emptycell & \emptycell & \emptycell & \emptycell & \emptycell & \emptycell \\
            \bottomrule
        \end{tabular}
    \end{subtable}
\end{table}

\subsection{Object erasure}
\label{sec:exp_object}

We evaluate ten ImageNet classes: church, golf ball, English springer, goose, chain saw, cassette player, French horn, garbage truck, banana, and parachute.
For each class, the prompt \texttt{a photo of \{class\}} is used with seeds 0--49, yielding 500 images per setting.
The single-target setting erases church, while the five-target setting erases church, golf ball, parachute, cassette player, and banana in the same model.
\Cref{tab:object_erasure} reports target Top-1 and Top-3 accuracy and non-target Top-1 accuracy, averaged over the corresponding classes.
The non-target average covers nine classes in the single-target setting and five in the five-target setting.

\begin{table}[t]
    \centering
    \small
    \setlength{\tabcolsep}{4pt}
    \caption{Object erasure over ten ImageNet classes. Accuracies (\%) are averaged over the target or non-target classes: one and nine classes in (a), and five and five in (b).}
    \label{tab:object_erasure}
    \begin{subtable}{\columnwidth}
        \centering
        \caption{Erasing church}
        \label{tab:object_single}
        \begin{tabular}{@{}lccc@{}}
            \toprule
            Method & \multicolumn{2}{c}{Target} & Non-target \\
            \cmidrule(lr){2-3}\cmidrule(lr){4-4}
            & Top-1 $\downarrow$ & Top-3 $\downarrow$ & Top-1 $\uparrow$ \\
            \midrule
            SD3.5 Medium & \emptycell & \emptycell & \emptycell \\
            UCE & \emptycell & \emptycell & \emptycell \\
            ESD & \emptycell & \emptycell & \emptycell \\
            EraseAnything & \emptycell & \emptycell & \emptycell \\
            RAZOR & \emptycell & \emptycell & \emptycell \\
            CASteer & \emptycell & \emptycell & \emptycell \\
            Ours & \emptycell & \emptycell & \emptycell \\
            \bottomrule
        \end{tabular}
    \end{subtable}
    \par\smallskip
    \begin{subtable}{\columnwidth}
        \centering
        \caption{Erasing five objects}
        \label{tab:object_multi}
        \begin{tabular}{@{}lccc@{}}
            \toprule
            Method & \multicolumn{2}{c}{Target} & Non-target \\
            \cmidrule(lr){2-3}\cmidrule(lr){4-4}
            & Top-1 $\downarrow$ & Top-3 $\downarrow$ & Top-1 $\uparrow$ \\
            \midrule
            SD3.5 Medium & \emptycell & \emptycell & \emptycell \\
            UCE & \emptycell & \emptycell & \emptycell \\
            ESD & \emptycell & \emptycell & \emptycell \\
            EraseAnything & \emptycell & \emptycell & \emptycell \\
            RAZOR & \emptycell & \emptycell & \emptycell \\
            CASteer & \emptycell & \emptycell & \emptycell \\
            Ours & \emptycell & \emptycell & \emptycell \\
            \bottomrule
        \end{tabular}
    \end{subtable}
\end{table}
\FloatBarrier
